\begin{remark}[Generalization of Reflex Fields Beyond CM Fields] \label{remark:generalization_of_reflex_field_beyond_cm_fields}
    The notion of \CrefAndHyperrefIfExist{definition:cm_type_and_reflex_field_of_a_cm_type}{reflex field of a CM type} generalizes to abelian varieties with other endomorphism types via the \CrefAndHyperrefIfExist{definition:rational_hodge_type_of_a_complex_abelian_variety}{rational Hodge type}. See \Cref{definition:reflex_field_of_the_rational_hodge_type_of_a_complex_abelian_variety}.

    This construction applies when $\End^0(A)$ is any semisimple $\mathbb{Q}$-algebra satisfying $[\End^0(A):\mathbb{Q}] \leq 2g$, including:
    \begin{itemize}
        \item $\End^0(A) = E$ a \CrefAndHyperrefIfExist{definition:cm_field}{CM field} (the classical case);
        \item $\End^0(A) = F$ a totally real field with $[F:\mathbb{Q}] = g$ (\CrefAndHyperrefIfExist{definition:abelian_variety_with_real_multiplication_over_a_field}{real multiplication});
        \item $\End^0(A)$ a quaternion algebra over a totally real field;
        \item $\End^0(A)$ an Albert type IV division algebra.
    \end{itemize}

    In the RM case, the Hodge type is unique up to conjugacy (there is only one choice of $\{1,0\}$-grading on each real embedding), and the reflex field equals the maximal totally real subfield.
\end{remark}
